\section{Modelos del mundo: de la fórmula y el control al cerebro del robot}

La parte técnica más sólida de este episodio es la definición de modelo del mundo que da Xie Saining (谢赛宁). No trata world model como una palabra de moda; regresa a las preguntas fundamentales del control, del model-based reinforcement learning y de la ciencia cognitiva: dados el estado actual y una acción, predecir el estado siguiente, y después usar esa predicción para guiar la planning y la decision making.

\[
s_{t+1} = f(s_t, a_t)
\]

Aquí, $s_t$ representa el estado del sistema o del entorno en el instante actual; $a_t$ representa la action o intervention que realiza el agente; $f$ es la transition / predictive function aprendida, y $s_{t+1}$ es el estado siguiente al que el sistema puede llegar después de ejecutar la acción. La dificultad real no está en la fórmula, sino en cómo obtener un $s_t$ adecuado a partir de percepciones de alta dimensión, cómo lograr que $f$ aprenda las leyes físicas y cómo usar la predicción para actuar en lugar de solo generar muestras vistosas.

\begin{figure}[H]
\centering
\includegraphics[width=0.96\textwidth]{world-model-stack.png}
\caption{El ciclo cerrado mínimo de un modelo del mundo: percepción, representación del estado, predicción de la transición, planificación y decisión, retroalimentación real y actualización del aprendizaje.}
\end{figure}

\begin{knowledgebox}{Lectura de la figura: el modelo del mundo no es un módulo aislado}
Lo más importante de la figura es el ciclo cerrado. Primero, la entrada perceptual se comprime en un state relevante para la tarea; el modelo predice el estado siguiente a partir del state y la action; el planificador compara distintas trayectorias futuras; al ejecutar la acción se obtiene retroalimentación real, y los errores y fallas regresan al entrenamiento. Si falta cualquiera de estos eslabones, el world model se reduce a un generador estático, a un sistema de preguntas y respuestas o a un modelo de representación fuera de línea.
\end{knowledgebox}

\subsection{La intuición del Model Predictive Control}

La fórmula anterior muestra la forma mínima de un world model; ahora hay que responder cómo se convierte en acción. El MPC ofrece la respuesta más sencilla y, a la vez, la de mayor valor didáctico: usar el modelo para imaginar el futuro y ejecutar únicamente el paso que más conviene en el momento actual.

En la entrevista se menciona el model predictive control. Su procedimiento básico es el siguiente: en el instante actual, se usa el modelo para hacer el roll out de varias action sequence futuras, se calcula el cost de cada trayectoria, se elige la secuencia de menor cost, se ejecuta su primer paso y en el instante siguiente se vuelve a planificar. No se trata de «expresar razones» al estilo de un LLM, sino de elegir acciones con base en la predicción de estados futuros.

\begin{knowledgebox}{Términos clave: conceptos relacionados con los modelos del mundo}
\begin{tabularx}{\textwidth}{p{0.23\textwidth}p{0.30\textwidth}X}
\toprule
Término & Explicación breve & Función en este episodio \\
\midrule
State & Información mínima suficiente que describe el estado actual del sistema & Vincula el aprendizaje de representaciones con la predicción de acciones. \\
Action / Intervention & Acción o intervención que el agente ejerce sobre el entorno & Hace que la predicción pase de la observación a la decisión. \\
Transition Function & Función que predice el state siguiente a partir del state y la action actuales & Forma central del modelo del mundo. \\
Planning & Comparar trayectorias futuras y elegir una acción & Pone la predicción del modelo al servicio de un objetivo. \\
MPC & Model Predictive Control, predicción y control con horizonte deslizante & Muestra cómo se usa un modelo del mundo para el control. \\
Model-based RL & Uso de un modelo del entorno como apoyo del aprendizaje por refuerzo & Explica la relación histórica entre el world model y el RL. \\
\bottomrule
\end{tabularx}
\end{knowledgebox}

\subsection{El world model es un propósito, no un algoritmo único}

Xie Saining subraya que el modelo del mundo es difícil de definir porque no es el nombre de un algoritmo, sino un propósito. Los modelos de lenguaje, los video diffusion model, las 3D representation, la robotics, los VLA y el model-based RL pueden avanzar hacia el world model desde direcciones distintas. Discutir «cuál es el verdadero world model» tiene sentido a corto plazo, pero a largo plazo la pregunta más importante es si el sistema puede comprender el mundo físico, conservar la memory relevante, reason and plan, hacer counterfactual / causal inference y ser controllable and safe.

\begin{importantbox}{Las cinco capacidades de un modelo del mundo}
Un world model orientado al mundo físico debe tener, como mínimo: physical world understanding, large associated memory, reasoning / planning, counterfactual or causal inference, y controllability and safety. No se mide con un indicador tan simple como «qué tan realista se ve el video generado».
\end{importantbox}

\subsection{Resumen de la sección}

La definición de fondo del modelo del mundo es sencilla: predecir el state que sigue a una action. La dificultad está en la representación del state, en el anclaje físico de la predicción, en la utilidad práctica de la planning y en el ciclo cerrado de retroalimentación. Tratar el world model como un objetivo, y no como una arquitectura de modelo concreta, evita que las disputas terminológicas de corto plazo desvíen la atención.
